\section{Strings}

\entry{KMP}{$O(n + m)$ \quad mem $n + m$}{%
From KACTL. \code{pi(s)[x]} is the length of the longest proper prefix of \code{s} that
ends at \code{x} (abacaba $\to$ 0010123). \code{match(s, pat)} returns every 0-indexed
start of \code{pat} in \code{s}, overlaps included; neither may contain
\code{'\textbackslash0'}.\\[0.3ex]
Shortest period of \code{s}: \code{n - pi(s).back()}; \code{s} is a repetition of it iff
it divides \code{n}. Borders of \code{s} are \code{n}, \code{p[n-1]},
\code{p[p[n-1]-1]}, \dots\ --- walk the chain. Counting occurrences of every prefix: add
1 at each \code{p} value and push counts down the chain in reverse order.}

%LABEL kmp
\begin{lstlisting}
vi pi(const string& s) {
    vi p(s.size());
    rep(i,1,(ll)s.size()) {
        int g = p[i-1];
        while (g && s[i] != s[g]) g = p[g-1];
        p[i] = g + (s[i] == s[g]);
    }
    return p;
}
vi match(const string& s, const string& pat) {
    vi p = pi(pat + '\0' + s), res;
    rep(i,(ll)p.size()-(ll)s.size(),(ll)p.size())
        if (p[i] == (ll)pat.size())
            res.push_back(i - 2 * (ll)pat.size());
    return res;
}
\end{lstlisting}

\entry{Z function}{$O(n)$ \quad mem $n$}{%
From KACTL. \code{Z(s)[i]} is the length of the longest common prefix of \code{s} and
\code{s[i..]}, with \code{z[0] = 0} (abacaba $\to$ 0010301). Pattern matching: build
\code{pat + '\textbackslash1' + txt} and look for \code{z[i] == pat.size()}; the
separator must not occur in either string. Often replaces KMP.}

%LABEL zfunc
\begin{lstlisting}
vi Z(const string& S) {
    vi z(S.size());
    int l = -1, r = -1;
    rep(i,1,(ll)S.size()) {
        z[i] = i >= r ? 0 : min<ll>(r - i, z[i - l]);
        while (i + z[i] < (ll)S.size()
               && S[i + z[i]] == S[z[i]])
            z[i]++;
        if (i + z[i] > r)
            l = i, r = i + z[i];
    }
    return z;
}
\end{lstlisting}

\entry{Manacher}{$O(n)$ \quad mem $n$}{%
From KACTL. Every palindromic substring at once: \code{auto p = manacher(s);}
\code{p[1][i]} = half the length (rounded down) of the longest \emph{odd} palindrome
centred at \code{i}; \code{p[0][i]} = half the length of the longest \emph{even} one
centred between \code{i-1} and \code{i}. Number of palindromic substrings:
$\sum (p[1][i] + 1) + \sum p[0][i]$.\\[0.3ex]
Is \code{s[l..r)} a palindrome, in $O(1)$: with \code{len = r-l}, \code{c = l+len/2},
it is \code{p[len\%2][c] >= len/2} --- e.g. for ``split into the fewest palindromes''
DP. Longest palindrome: $\max(2p[1][i]+1, 2p[0][i])$.}

%LABEL manacher
\begin{lstlisting}
array<vi, 2> manacher(const string& s) {
    int n = (int)s.size();
    array<vi,2> p = {vi(n+1), vi(n)};
    rep(z,0,2) for (int i=0,l=0,r=0; i < n; i++) {
        int t = r-i+!z;
        if (i<r) p[z][i] = min<ll>(t, p[z][l+t]);
        int L = i-p[z][i], R = i+p[z][i]-!z;
        while (L>=1 && R+1<n && s[L-1] == s[R+1])
            p[z][i]++, L--, R++;
        if (R>r) l=L, r=R;
    }
    return p;
}
\end{lstlisting}

\entry{String hashing}{build $O(n)$, compare $O(1)$ \quad mem $n$}{%
\code{get(l, r)} is the hash of \code{s[l..r)}; two substrings are equal iff their
hashes match. Two independent moduli, so a collision needs both to fail at once ---
with a single 64-bit overflow hash an adversarial test \emph{will} break you.\\[0.3ex]
This is the workhorse: compare any two substrings, binary-search the longest common
prefix of two suffixes, dedupe substrings by putting \code{get} results in a
\code{set}, Rabin-Karp matching. Often replaces a suffix array
outright.\\[0.3ex]
For randomised bases (safer against hash-killers) pick \code{B1}, \code{B2} at run time
from \code{mt19937} in $[256, M)$.}

%LABEL hash
\begin{lstlisting}
struct Hash {
    static constexpr ll M1 = 1000000007, M2 = 998244353,
                        B1 = 131, B2 = 137;
    int n; vi h1, h2, p1, p2;
    Hash(const string& s) : n((int)s.size()), h1(n+1, 0),
            h2(n+1, 0), p1(n+1, 1), p2(n+1, 1) {
        rep(i, 0, n) {
            h1[i+1] = (h1[i] * B1 + s[i]) % M1;
            h2[i+1] = (h2[i] * B2 + s[i]) % M2;
            p1[i+1] = p1[i] * B1 % M1;
            p2[i+1] = p2[i] * B2 % M2;
        }
    }
    pii get(int l, int r) {                   // s[l..r)
        ll a = (h1[r] - h1[l] * p1[r-l]) % M1;
        ll b = (h2[r] - h2[l] * p2[r-l]) % M2;
        return {(a + M1) % M1, (b + M2) % M2};
    }
};
\end{lstlisting}

\entry{Suffix array + LCP}{$O(n \log n)$ \quad mem $n$}{%
From KACTL. \code{SuffixArray S(s);} \code{S.sa} has size $n+1$: \code{sa[0] = n} is
the empty suffix, then \code{sa[i]} is the start of the \code{i}-th smallest suffix.
\code{lcp[i]} = longest common prefix of the suffixes \code{sa[i]} and \code{sa[i-1]}
(\code{lcp[0] = 0}). \code{s} must not contain \code{'\textbackslash0'}. For a
\code{vector<int>} input use values $\ge 1$ and set \code{lim} above the
largest.\\[0.3ex]
Number of distinct substrings is $\frac{n(n+1)}{2} - \sum \mathrm{lcp}$. Longest
repeated substring is $\max \mathrm{lcp}$. Longest common substring of two strings: join
with a separator and take the best \code{lcp} between neighbours from different sides.
LCP of two arbitrary suffixes = min of \code{lcp} between their ranks --- feed
\code{lcp} to \code{RMQ} (rank = inverse of \code{sa}).\\[0.3ex]
Consider \code{Hash} plus binary search first: shorter to write, usually fast enough.}

%LABEL sufarr
\begin{lstlisting}
struct SuffixArray {
    vi sa, lcp;
    SuffixArray(string s, int lim=256) { // or vector<int>
        s.push_back(0); int n = (int)s.size(), k = 0, a, b;
        vi x(all(s)), y(n), ws(max(n, lim));
        sa = lcp = y, iota(all(sa), 0);
        for (int j = 0, p = 0; p < n;
             j = max(1, j * 2), lim = p) {
            p = j, iota(all(y), n - j);
            rep(i,0,n) if (sa[i] >= j) y[p++] = sa[i] - j;
            fill(all(ws), 0);
            rep(i,0,n) ws[x[i]]++;
            rep(i,1,lim) ws[i] += ws[i - 1];
            for (int i = n; i--;) sa[--ws[x[y[i]]]] = y[i];
            swap(x, y), p = 1, x[sa[0]] = 0;
            rep(i,1,n) a = sa[i - 1], b = sa[i], x[b] =
                (y[a] == y[b] && y[a + j] == y[b + j])
                    ? p - 1 : p++;
        }
        for (int i = 0, j; i < n - 1; lcp[x[i++]] = k)
            for (k && k--, j = sa[x[i] - 1];
                    s[i + k] == s[j + k]; k++);
    }
};
\end{lstlisting}

\entry{Aho-Corasick}{build $O(26 N)$, \code{find} $O(n)$ \quad mem $26 N$}{%
From KACTL. Many patterns at once. \code{AhoCorasick ac(pat);} ($N$ = total pattern
length, no empty patterns). \code{ac.find(word)[i]} = index of the
\textbf{longest} pattern ending at position \code{i}, or $-1$.
\code{ac.findAll(pat, word)[i]} = all patterns \textbf{starting} at \code{i} --- that
can be $N\sqrt N$ matches, so only when the output is small. Alphabet: set
\code{alpha} and \code{first} (\code{'a'} for lowercase).\\[0.3ex]
\textbf{Count} matches instead of listing them: \code{nmatches} of a node is the number
of patterns ending there, fail links included --- uncomment the \code{count} lines in
\code{find}, $O(n)$ total. Duplicate patterns are fine; \code{backp} chains them.\\[0.3ex]
DP over ``text avoiding all patterns'': states are node ids, transitions
\code{N[n].next[c]}, forbidden where \code{N[n].end != -1}.}

%LABEL aho
\begin{lstlisting}
struct AhoCorasick {
    enum {alpha = 26, first = 'a'}; // change this!
    struct Node {
        // (nmatches is optional)
        int back, next[alpha], start = -1, end = -1,
            nmatches = 0;
        Node(int v) { memset(next, v, sizeof(next)); }
    };
    vector<Node> N;
    vector<int> backp;     // int: assigned via ?: below
    void insert(string& s, int j) {
        assert(!s.empty());
        int n = 0;
        for (char c : s) {
            int& m = N[n].next[c - first];
            if (m == -1) {
                n = m = (int)N.size(); N.emplace_back(-1);
            }
            else n = m;
        }
        if (N[n].end == -1) N[n].start = j;
        backp.push_back(N[n].end);
        N[n].end = j;
        N[n].nmatches++;
    }
    AhoCorasick(vector<string>& pat) : N(1, -1) {
        rep(i,0,(ll)pat.size()) insert(pat[i], i);
        N[0].back = (int)N.size();
        N.emplace_back(0);
        queue<int> q;
        for (q.push(0); !q.empty(); q.pop()) {
            int n = q.front(), prev = N[n].back;
            rep(i,0,alpha) {
                int &ed = N[n].next[i], y = N[prev].next[i];
                if (ed == -1) ed = y;
                else {
                    N[ed].back = y;
                    (N[ed].end == -1 ? N[ed].end
                        : backp[N[ed].start]) = N[y].end;
                    N[ed].nmatches += N[y].nmatches;
                    q.push(ed);
                }
            }
        }
    }
    vi find(string word) {
        int n = 0;
        vi res; // ll count = 0;
        for (char c : word) {
            n = N[n].next[c - first];
            res.push_back(N[n].end);
            // count += N[n].nmatches;
        }
        return res;
    }
    vector<vi> findAll(vector<string>& pat, string word) {
        vi r = find(word);
        vector<vi> res(word.size());
        rep(i,0,(ll)word.size()) {
            int ind = r[i];
            while (ind != -1) {
                res[i - pat[ind].size() + 1].push_back(ind);
                ind = backp[ind];
            }
        }
        return res;
    }
};
\end{lstlisting}

\entry{XOR trie}{insert / query $O(\mathrm{BITS})$ \quad mem $n \cdot \mathrm{BITS}$}{%
\code{maxXor(x)} is the largest \code{x \textasciicircum{} y} over everything inserted:
walk from the top bit and always take the opposite branch if it exists. Raise
\code{BITS} to 62 for \code{ll} values.\\[0.3ex]
Keep a counter per node and you can also answer ``how many \code{y} with
\code{x \textasciicircum{} y < k}'', and delete by decrementing. Maximum XOR
\emph{subarray}: insert prefix xors as you sweep.\\[0.3ex]
A plain prefix trie over letters is the same shape with 26 children --- see
\code{AhoCorasick} without the fail links.}

\begin{lstlisting}
struct XorTrie {
    static const int BITS = 30;
    vector<array<int,2>> ch{{-1, -1}};
    void insert(ll x) {
        int u = 0;
        down(b, BITS + 1, 0) {
            int c = x >> b & 1;
            if (ch[u][c] < 0) {
                ch.push_back({-1, -1});
                ch[u][c] = (int)ch.size() - 1;
            }
            u = ch[u][c];
        }
    }
    ll maxXor(ll x) {         // needs >= 1 inserted value
        int u = 0; ll r = 0;
        down(b, BITS + 1, 0) {
            int c = (x >> b & 1) ^ 1;
            if (ch[u][c] >= 0) { r |= 1LL<<b; u = ch[u][c]; }
            else u = ch[u][c ^ 1];
        }
        return r;
    }
};
\end{lstlisting}

\entry{Minimum rotation}{$O(n)$}{%
Index where the lexicographically smallest rotation of \code{s} starts:
\code{rotate(s.begin(), s.begin() + minRotation(s), s.end())}. Two strings (or
necklaces, or cyclic sequences of numbers) are rotations of each other iff their
minimum rotations are equal --- the standard way to put a cycle into one canonical form
before hashing or comparing. Works on \code{vector<int>} too if you change the type.
From KACTL.}

%LABEL minrot
\begin{lstlisting}
int minRotation(string s) {
    int a = 0, N = (int)s.size(); s += s;
    rep(b, 0, N) rep(k, 0, N) {
        if (a + k == b || s[a+k] < s[b+k]) {
            b += max(0, (int)k - 1); break;
        }
        if (s[a+k] > s[b+k]) { a = (int)b; break; }
    }
    return a;
}
\end{lstlisting}
